% ============================================================================
% CAPÍTULO 3: DPI FORMAL — LA DESIGUALDAD DE PROCESAMIENTO DE DATOS
% ============================================================================
\chapter{Fundamentos de Teoría de la Información: La DPI y la Destrucción Irreversible}
\label{ch:dpi}

\epigraph{``La entropía siempre crece.''\\
Paráfrasis del Segundo Principio de la Termodinámica.\\[0.3cm]
``La información semántica nunca aumenta al tokenizar.''\\
Consecuencia directa de la DPI (Shannon, 1948).}{--- Axiomas del dogma No-Worm}

\section{Preliminares: Teoría de la Información de Shannon}

\subsection{Entropía de Shannon}

Sea $X$ una variable aleatoria discreta con distribución de probabilidad
$\{p_1, p_2, \ldots, p_n\}$. La entropía de Shannon de $X$ es:

\begin{equation}
H(X) = -\sum_{i=1}^{n} p_i \log_2 p_i \quad \text{(bits)}
\label{eq:shannon_entropy}
\end{equation}

La entropía cuantifica la \textit{información promedio} contenida en una realización
de $X$. Para una variable continua $X \in \R^D$ con densidad $f(\mathbf{x})$:

\begin{equation}
h(X) = -\int_{\R^D} f(\mathbf{x}) \ln f(\mathbf{x}) \, d\mathbf{x}
\label{eq:differential_entropy}
\end{equation}

Para $\mathbf{h} \sim \mathcal{N}(\mathbf{0}, \sigma^2 I_D)$ (aproximación al
estado latente de un LLM con temperatura alta):

\begin{equation}
h(\mathbf{h}) = \frac{D}{2} \ln(2\pi e \sigma^2)
\end{equation}

Este valor crece linealmente con $D$, confirmando que la capacidad informacional
del espacio latente escala con la dimensionalidad.

\subsection{Información Mutua}

La información mutua entre $X$ e $Y$ mide la dependencia estadística:

\begin{equation}
I(X; Y) = H(X) - H(X|Y) = H(Y) - H(Y|X) = H(X) + H(Y) - H(X,Y)
\label{eq:mutual_info}
\end{equation}

\begin{theorem}[Propiedades de la Información Mutua]
\label{thm:mi_props}
\begin{enumerate}
\item $I(X; Y) \ge 0$, con igualdad si y solo si $X \perp Y$.
\item $I(X; Y) \le \min\{H(X), H(Y)\}$ (la información mutua no puede exceder
la entropía de ninguna de las variables).
\item $I(X; Y) = I(Y; X)$ (simetría).
\end{enumerate}
\end{theorem}

\section{La Data Processing Inequality (DPI)}

\begin{theorem}[Data Processing Inequality]
\label{thm:dpi_formal}
Sea $(X, Y, Z)$ una cadena de Markov, es decir, $X \to Y \to Z$
(lo que significa que $Z$ es condicionalmente independiente de $X$ dado $Y$:
$P(Z|X,Y) = P(Z|Y)$). Entonces:
\begin{equation}
I(X; Z) \le I(X; Y)
\label{eq:dpi}
\end{equation}
y la igualdad se alcanza si y solo si la función $Y \to Z$ es suficiente
para $X$, es decir, si $(X, Z)$ también forman una cadena de Markov.
\end{theorem}

\begin{proof}
Por la regla de la cadena de la información mutua:
\begin{align}
I(X; Y, Z) &= I(X; Y) + I(X; Z | Y) \\
I(X; Y, Z) &= I(X; Z) + I(X; Y | Z)
\end{align}
Como $X \to Y \to Z$ es una cadena de Markov, $X$ es condicionalmente independiente
de $Z$ dado $Y$, por lo tanto $I(X; Z | Y) = 0$. Así:
\begin{equation}
I(X; Y) = I(X; Y, Z) = I(X; Z) + I(X; Y | Z) \ge I(X; Z)
\end{equation}
donde la desigualdad final usa que $I(X; Y | Z) \ge 0$.
\end{proof}

\section{Aplicación al Pipeline de Tokenización}

\subsection{Formalización de la Cadena de Markov}

Definimos las variables aleatorias:
\begin{itemize}
\item $Q$ = consulta del usuario (query original)
\item $\mathbf{h} \in \R^D$ = estado oculto final del LLM dado $Q$
\item $\mathbf{l} \in \R^{|V|}$ = logits (proyección de $\mathbf{h}$ mediante $W_E$)
\item $\mathbf{p} = \text{Softmax}(\mathbf{l}/T) \in \Delta^{|V|-1}$ = distribución de tokens
\item $t \in \{0,\ldots,|V|-1\}$ = token muestreado
\item $\mathbf{e}_t \in \R^{D'}$ = embedding del token en el agente receptor
\end{itemize}

El pipeline completo forma la cadena de Markov:
\begin{equation}
Q \to \mathbf{h} \to \mathbf{l} \to \mathbf{p} \to t \to \mathbf{e}_t
\end{equation}

Por aplicaciones sucesivas del Teorema~\ref{thm:dpi_formal}:
\begin{equation}
I(Q; \mathbf{e}_t) \le I(Q; t) \le I(Q; \mathbf{p}) \le I(Q; \mathbf{l}) \le I(Q; \mathbf{h})
\label{eq:dpi_chain}
\end{equation}

\begin{corollary}[Pérdida Estricta de Información por Cuantización Finitaria]
\label{cor:dpi_acumulada}
En el pipeline de tokenización $Q \to \mathbf{h} \to \mathbf{l} \to \mathbf{p} \to t \to \mathbf{e}_t$,
la información sobre el estado latente $\mathbf{h}$ o la consulta $Q$ disponible
en el embedding del receptor $\mathbf{e}_t$ satisface:
\begin{equation}
I(Q; \mathbf{e}_t) \le I(Q; t) \le I(Q; \mathbf{h}) - H(\mathbf{h} \mid t)
\end{equation}
donde la desigualdad es \textbf{estricta} ($I(Q; t) < I(Q; \mathbf{h})$) siempre que 
el soporte efectivo de la distribución de $\mathbf{h}$ tenga cardinalidad 
superior al tamaño del vocabulario: $|\text{supp}(\mathbf{h})| > |V|$.
\end{corollary}

\begin{proof}
La aplicación de muestreo o cuantización $\mathbf{h} \mapsto t \in \{0, \dots, |V|-1\}$ 
particiona el espacio continuo $\R^D$ (o la esfera $S^{D-1}$) en $|V|$ regiones de Voronoi 
$\{R_k\}_{k=1}^{|V|}$. Como $|\text{supp}(\mathbf{h})| > |V|$, existe al menos una región $R_k$ 
con medida no nula que contiene infinitos estados $\mathbf{h}_1 \neq \mathbf{h}_2$ mapeados al 
mismo token $t = k$. La aplicación es estrictamente no-inyectiva. Por el teorema de suficiencia 
de Kolmogorov-Blackwell, $t$ es un estadístico suficiente para $\mathbf{h}$ si y solo si la 
distribución condicional $P(\mathbf{h} \mid t)$ no depende de $Q$. Dado que $\mathbf{h}$ preserva 
grados de libertad continuos ortogonales a la proyección de logits, $H(\mathbf{h} \mid t) > 0$, 
lo que garantiza $\Delta I = I(Q; \mathbf{h}) - I(Q; t) > 0$. \qed
\end{proof}

\subsection{Cuantificación de la Pérdida}

Bajo la hipótesis de que el estado latente $\mathbf{h}$ sigue una distribución
gaussiana isotrópica en el subespacio de la consulta:

\begin{align}
I(Q; \mathbf{h}) &\approx \frac{D}{2} \log_2\left(1 + \text{SNR}_{\mathbf{h}}\right) \quad \text{(bits)} \\
I(Q; t) &\le \log_2 |V| \approx 15 \text{ bits} \quad (|V| = 32768)
\end{align}

Para $D = 4096$ y $\text{SNR}_{\mathbf{h}} \approx 10$ (valor típico en transformers):
\begin{equation}
I(Q; \mathbf{h}) \approx \frac{4096}{2} \log_2(11) \approx 17240 \text{ bits}
\end{equation}

La pérdida absoluta por tokenización es:
\begin{equation}
\Delta I = I(Q; \mathbf{h}) - I(Q; t) \ge 17240 - 15 = 17225 \text{ bits} \approx 2153 \text{ bytes}
\end{equation}

Esto equivale al $99.91\%$ de la información disponible.

\section{La DPI y el Enjambre Multi-Agente}

En un sistema de $N$ agentes donde agentes $i$ y $j$ se comunican mediante
tokens, la información disponible en el agente $j$ sobre el estado del agente $i$
satisface:

\begin{equation}
I(\mathbf{h}_i; \mathbf{h}_j^{\text{actualizado}}) \le I(\mathbf{h}_i; t_{ij}) \le I(\mathbf{h}_i; \mathbf{h}_i) = H(\mathbf{h}_i)
\end{equation}

Con $E$ aristas de comunicación en el enjambre, la pérdida total por ronda es:

\begin{equation}
\Delta I_{\text{total}} = \sum_{(i,j) \in E} \left[I(\mathbf{h}_i; \mathbf{h}_i) - I(\mathbf{h}_i; t_{ij})\right]
\ge E \cdot (D \cdot \log_2 \sigma_{\min} - \log_2 |V|)
\end{equation}

Para el grafo completo $K_{10}$ (enjambre de 10 agentes) con $E = 45$ aristas:
\begin{equation}
\Delta I_{K_{10}} \ge 45 \cdot (D \cdot 3.32 - 15) \text{ bits}
\end{equation}

Con $D = 8192$: pérdida de $\ge 45 \times 27207 \approx 1.22 \times 10^6$ bits por ronda.
Esto es \textbf{152 KB de información semántica destruida por ronda de comunicación}.

\section{El Principio de No-Aumento de la Entropía Semántica}

\begin{theorem}[No-Aumento de Entropía Semántica]
\label{thm:no_aumento}
Sea $\mathcal{S}$ el contenido semántico de un estado latente $\mathbf{h} \in \R^D$,
definido como la clase de equivalencia de todos los estados que darían la misma
respuesta de comportamiento. Entonces para cualquier pipeline de procesamiento
que involucre tokenización intermedia:

\begin{equation}
I(\mathcal{S}^{(t+1)}; Q) \le I(\mathcal{S}^{(t)}; Q)
\end{equation}

con igualdad si y solo si el canal de comunicación es suficiente, lo que requiere
que la capacidad del canal satisfaga $C \ge H(\mathcal{S} | Q)$.
\end{theorem}

\begin{proof}
El resultado es consecuencia directa de la DPI aplicada a la cadena de Markov
$Q \to \mathbf{h}^{(t)} \to t^{(t)} \to \mathbf{h}^{(t+1)}$. La suficiencia
del canal requiere $C = \log_2|V| \ge H(\mathbf{h}|Q)$, lo que es imposible
para $D \gg \log_2|V|$.
\end{proof}

\section{POLYDIM como Canal Suficiente}

El Protocolo PMTP propone un canal que evita la tokenización intermedia:

\begin{equation}
\mathbf{h}_A \xrightarrow{\text{PMTP}} \mathbf{h}_A' \in \R^{D}
\end{equation}

donde $\mathbf{h}_A' = \mathbf{h}_A$ exactamente (copia bit-a-bit), o
$\mathbf{h}_A' = M \cdot \mathbf{h}_A$ donde $M \in O(D')$ es una isometría
que alinea los espacios latentes heterogéneos.

En el caso de copia exacta (mismo espacio latente):
\begin{equation}
I(\mathbf{h}_A; \mathbf{h}_A') = H(\mathbf{h}_A) = I(\mathbf{h}_A; Q)
\end{equation}

Es decir, PMTP alcanza el límite superior de la DPI: el canal es suficiente.
La Información Mutua se preserva al $100\%$.

\begin{certifiedbox}
\textbf{Verificación empírica (V762, Exit Code 0):}
El PMTP Zero-Copy fue verificado con $D = 10^6$, transferencia de $8\,\text{MB}$.
Resultado: Max Bit Difference $= 0.0000 \times 10^{+0}$.
La copia bit-a-bit exacta confirma el canal suficiente.
\end{certifiedbox}

\section{La Frontera de Huffman y los Límites del Protocolo PMTP}

Una pregunta natural: ¿podría una compresión sin pérdida (Huffman, Lempel-Ziv)
preservar la información mientras reduce el ancho de banda del canal?

La respuesta es negativa para el caso de uso de POLYDIM por tres razones:

\begin{enumerate}
\item \textbf{Overhead de compresión:} Para vectores aleatorios en $\R^D$ con
distribución uniforme (como los embeddings normalizados en $S^{D-1}$), la entropía
de Shannon alcanza el máximo y la compresión sin pérdida no reduce el tamaño.

\item \textbf{Latencia de compresión:} Los algoritmos de compresión sin pérdida
tienen complejidad $\Omega(D)$ con constantes grandes. Para $D = 10^6$, la
latencia de compresión excede la latencia de transferencia en red local.

\item \textbf{Restricción topológica:} El espacio $S^{D-1}$ no tiene estructura
algebraica que permita compresión sin pérdida eficiente: no existe una base
``natural'' en la que la mayor parte de los coeficientes sean cero.
\end{enumerate}

La única compresión que no destruye información es la compresión en el espacio
tangente $T_{\mathbf{h}} S^{D-1}$ mediante proyecciones isométricas (Stiefel),
que es exactamente lo que implementa la retracción Cayley-SMW estudiada en el
Capítulo~\ref{ch:cayley_smw}.

\section{Entropía Topológica y el Número de Betti como Preservador de Semántica}

La teoría de la información es insuficiente para capturar todos los aspectos
de la semántica latente. Complementariamente, la topología algebraica ofrece
invariantes que miden la estructura ``forma'' del espacio semántico:

\begin{definition}[Número de Betti en Contexto POLYDIM]
\label{def:betti}
Para un espacio de estados semánticos $\mathcal{X}$ modelado como un complejo
simplicial, el $k$-ésimo número de Betti $\beta_k$ es el rango del $k$-ésimo
grupo de homología $H_k(\mathcal{X}; \mathbb{Z})$:
\begin{itemize}
\item $\beta_0$: número de componentes conexas (agentes desconectados del enjambre).
\item $\beta_1$: número de ``agujeros'' topológicos 1D (rutas alternativas de consenso).
\item $\beta_2$: número de ``cavidades'' 2D (subespacios de divergencia semántica).
\end{itemize}
\end{definition}

El \textbf{Betti-0 Guard} de POLYDIM verifica en tiempo real que $\beta_0 = 1$
(enjambre conexo) y $\beta_1 > 0$ (al menos un ciclo alternativo, garantizando
redundancia). La destrucción semántica por tokenización aumenta $\beta_0$ (el
enjambre se fragmenta) y puede decrementar $\beta_1$ (pérdida de caminos
alternativos de consenso).

\begin{certifiedbox}
\textbf{Betti-1 Guard V762:}
\texttt{Connected = 0 (PASS: $\beta_0 = 1$)},
\texttt{Fragmented = -6 (PASS: topología intacta)}.
El enjambre de 8 agentes mantuvo $\beta_0 = 1$ durante todas las pruebas.
\end{certifiedbox}
